%======================================================================
\section{Related Work}
\label{sec:related}
%======================================================================

\paragraph{Multi-Agent LLM Consensus.}
The aggregation of multiple LLM outputs has been studied through several paradigms. Self-consistency~\cite{wang2023selfconsistency} samples multiple reasoning paths and selects the most frequent answer. Multi-agent debate~\cite{du2024improving,choi2025debate} extends this to iterative discussion among multiple model instances. ChatEval~\cite{chan2023chateval} applies this to evaluation tasks. Zhao et al.~\cite{zhao2024electoral} propose cardinal voting based on social choice theory. Li et al.~\cite{li2025dynamic} introduce dynamic weighting.

A critical failure mode in these systems is \emph{sycophancy}---agents adopting majority positions without genuine reasoning. Bertalani\v{c} and Fortuna~\cite{bertalanic2026cost} provide the most comprehensive empirical characterization to date, documenting three failure pathways in homogeneous multi-agent debate ($n{=}10$): sycophantic conformity (up to 85.5\% modal adoption), contextual fragility (up to 70\% of previously correct agents destabilized by peer rationales), and consensus collapse (up to 32.3pp oracle gap where plurality voting discards correct answers). Their central finding---that homogeneous teams without structured roles do not benefit from unguided peer exchange---directly motivates our work on structured role enforcement through AEIC decomposition and behavioral regulation through Stackelberg mechanisms.

Pitre et al.~\cite{pitre2025consensagent} address sycophancy through ConsensAgent, which dynamically refines prompts based on agent interactions to mitigate herding. ConsensAgent operates at the prompt level within a single debate topology. Our dynamic regulation differs fundamentally: rather than optimizing prompts for individual agents, we monitor \emph{collective behavioral dynamics} (stance change rate, majority ratio, convergence speed) across rounds and apply differentiated interventions---parameter adjustment \emph{and} prompt modification---targeted at the specific pathology detected (sycophantic vs.\ deadlock). Swamy et al.~\cite{swamy2025peacemaker} provide formal metrics for sycophancy in multi-agent debate; their finding that sycophancy levels vary substantially across agent roles informs our stance-change-based detection. Most recently, Kumarappan et al.~\cite{kumarappan2026alignment} demonstrate that multi-agent sycophancy is structural (present in base models, not just RLHF artifacts) and that mitigations should target ``structured dissent at the pipeline level''---this is precisely what our Stackelberg dynamic regulation achieves through behavioral state detection and adaptive EV-weight suppression.

\paragraph{Structured Reasoning Decomposition.}
Chain-of-Thought prompting~\cite{wei2022chain} elicits step-by-step reasoning, and Tree-of-Thoughts~\cite{yao2024tree} extends this to structured exploration of reasoning paths. ReAct~\cite{yao2023react} interleaves reasoning with action. These methods decompose reasoning to improve \emph{individual} answer quality. Our AEIC decomposition serves a different purpose: it structures agent outputs into four parseable fields (Assumption, Evidence, Inference, Conclusion) to enable \emph{cross-agent} operations---reliable stance extraction from the explicit Conclusion~field (+28pp over unstructured extraction), cross-topology voting-pattern comparison for adaptive routing, and population-level behavioral dynamics monitoring for the dynamic regulation mechanism (Section~\ref{sec:dynamic}). The decomposition is not designed to improve per-agent reasoning; it is designed to make collective behavior observable and regulatable.

\paragraph{Communication Topology in Multi-Agent Systems.}
The effect of communication structure on multi-agent performance has been explored in several contexts. Chen et al.~\cite{chen2023agentverse} use a centralized evaluator (star topology). Sequential architectures appear in chain-of-agents designs~\cite{hong2023metagpt}. Debate topologies (complete graphs) are studied by Du et al.~\cite{du2024improving}. Chen et al.~\cite{chen2024reconcile} introduce round-table conference among diverse LLMs with confidence-weighted voting. Most recently, RUMAD~\cite{wang2026rumad} applies reinforcement learning to adaptively control debate rounds and agent selection within a single topology type. Our approach differs in three ways: we select among five fundamentally different communication structures (parallel, sequential, star, debate, hierarchical), use interpretable voting-pattern features rather than opaque RL policies, and ground our mechanism in game-theoretic incentive compatibility guarantees.

\paragraph{Game Theory and LLMs.}
Sun et al.~\cite{sun2025gametheory} provide a taxonomy of game theory applied to LLMs. Yi et al.~\cite{yi2025debate} model multi-agent debate as a Bayesian game with Nash equilibria. Rauba et al.~\cite{rauba2025principal} argue for principal--agent modeling of LLM systems. Liu et al.~\cite{liu2025interpretable} use LLMs for automated mechanism design. Our Stackelberg formulation differs from simultaneous-move equilibrium approaches by modeling the mechanism designer as a leader who \emph{commits} to aggregation rules before agents act---a more natural fit for mechanism design where the goal is shaping follower behavior through credible commitment.

\paragraph{Stackelberg Games in Distributed Systems.}
Stackelberg incentive mechanisms are established in federated learning~\cite{fu2024stackelberg,guo2023incentive} and crowdsourcing~\cite{wang2024incentive}. Saig et al.~\cite{saig2024incentivizing} apply contract theory to incentivize quality text generation from LLMs. Xu et al.~\cite{xu2026rcbsf} apply Stackelberg games to multi-agent contract revision. To the best of our knowledge, no prior work applies Stackelberg mechanism design to multi-agent LLM consensus with adaptive topology routing.

\paragraph{Positioning against Closest Work.}
Table~\ref{tab:related} positions our framework against the four most closely related approaches.

\begin{table}[t]
\caption{Comparison with closest related work.}\label{tab:related}
\centering\small
\begin{tabular}{@{}lccccc@{}}
\toprule
& \makebox[1.8cm]{\shortstack{Formal\\Incentives}} & \makebox[1.6cm]{\shortstack{Topology\\Selection}} & \makebox[1.6cm]{\shortstack{Sycophancy\\Solution}} & \makebox[1.5cm]{\shortstack{Structured\\Output}} & \makebox[1.5cm]{\shortstack{Behavioral\\Regulation}} \\
\midrule
Standard Debate~\cite{du2024improving} & --- & --- & --- & --- & --- \\
ConsensAgent~\cite{pitre2025consensagent} & --- & --- & Prompt-level & --- & --- \\
RUMAD~\cite{wang2026rumad} & --- & \checkmark$^\dagger$ & --- & --- & --- \\
Bertalani\v{c} et al.~\cite{bertalanic2026cost} & --- & --- & Analysis only & --- & --- \\
\midrule
\textbf{This work} & \checkmark$^\ddagger$ & \checkmark & \checkmark & \checkmark & \checkmark \\
\bottomrule
\end{tabular}
\\[2pt]
{\footnotesize $^\dagger$Within-topology only (varying debate rounds/participation). $^\ddagger$Stackelberg $\epsilon$-incentive compatibility (Theorem~\ref{thm:ic}).}
\end{table}
